\documentclass[10pt,a4paper]{amsart}
\usepackage[margin=19mm]{geometry}
\usepackage{amsmath,amssymb,mathtools}
\usepackage[scr=rsfs,cal=euler]{mathalfa}
\usepackage{xfrac}
\usepackage[unicode,hidelinks,bookmarks=false]{hyperref}
\newcommand{\ie}{i.e.\ }
\newcommand{\cf}{cf.\ }
\newcommand{\resp}{respectively\ }
\newcommand{\Subsection}{\subsection}
\providecommand{\nobreakdash}{\nobreak\hbox{-}}
\setlength{\emergencystretch}{2em}
\multlinegap0pt
\usepackage{bm}
\usepackage{bbm}
\usepackage{dsfont}
\usepackage{xspace}
\usepackage{cancel}
\usepackage{mathtools}
\usepackage{amsthm}
\makeatletter
\@ifundefined{theorem}{\newtheorem{theorem}{Theorem}}{}
\@ifundefined{definition}{\newtheorem{definition}[theorem]{Definition}}{}
\@ifundefined{proposition}{\newtheorem{proposition}{Proposition}}{}
\makeatother
\def\C{\mathbb C}
\def\DS{\mathcal{D}_S}
\def\R{\mathbb R}
\def\calB{\mathcal B}
\def\calD{\mathcal D}
\def\im{\text{im}}
\def\tr{\text{tr}}
\title[Best approximants relative to a C$^*$-subalgebra, joint nume]{Best approximants relative to a C$^*$-subalgebra, joint numerical range and subdifferentials}
\author{Tamara Bottazzi, Alejandro Varela}
\date{Source: arXiv:2604.16041v1 (CC BY 4.0)}
\begin{document}
\maketitle

\noindent\emph{Source and licence.} Best approximants relative to a C$^*$-subalgebra, joint numerical range and subdifferentials. Version arXiv:2604.16041v1 (CC BY 4.0), distributed under the Creative Commons Attribution 4.0 International licence, as recorded in the arXiv licence metadata for this exact version and shown on the abstract page https://arxiv.org/abs/2604.16041v1. Full rights notice: licenses/arxiv-2604.16041-CC-BY-4.0.txt.\par\smallskip

\noindent\textit{Editorial scope.} Counted unit: the Proposition whose source label is \texttt{prop 1} in the article (source0.tex lines 367--411, source label \texttt{prop 1}), delivered together with the article's own complete original proof of it (source0.tex lines 412--526, one \texttt{proof} environment of 115 lines). No mathematics is written, repaired or re-derived: statement and proof are the article's own text line for line. Required context retained uncounted: def jnr (lines 251--255); defi mS (lines 257--268); def momento de un Subesp en relacion a B (lines 259--264); def matrices de densidad de S (lines 271--274); equ densidad escrita usando autovals y autovecs (lines 278--280); defi phi (lines 284--286); wi (lines 354--356). Every definition, display and label of those lines is retained unchanged. Omitted from this package and not redistributed: 1--250, 256--256, 265--270, 275--277, 281--283, 287--353, 357--366, 527--1184. Nothing inside a retained excerpt is omitted or altered: every retained body is delivered exactly as the article's lines are written, so no omission receipt of any kind accompanies this package. Citations: the retained excerpts carry no citation command at all, so no bibliography entry of the article is transcribed and no reference is redistributed. Reference adaptation: none was needed and none was made; every reference inside the delivered unit resolves inside it, so no unresolved reference or citation is produced. Printed slips kept literal: no printed word, symbol, hypothesis or number of the counted statement or of its proof was altered. Layout and class: the article's own class is \texttt{article} with packages including authblk, inputenc, babel, amsfonts, amsmath, mathabx, fancyhdr, graphicx, xcolor, caption, ulem, mathtools, enumitem, amsthm; the wrapper is the lane's pinned \texttt{amsart} editorial wrapper at 10pt on A4, which restates the theorem-like environments the retained text needs and adds the article's own macro definitions that the retained text needs (\texttt{\textbackslash C}, \texttt{\textbackslash DS}, \texttt{\textbackslash R}, \texttt{\textbackslash calB}, \texttt{\textbackslash calD}, \texttt{\textbackslash im}, \texttt{\textbackslash tr}), with extra wrapper packages (bm, bbm, dsfont, xspace, cancel, mathtools). The delivered numbering is the unit's own. Assets: no retained span contains a figure environment, a graphics-inclusion command, a PDF-page inclusion, a table or a TikZ picture; no image file is copied into this package, the package declares no asset dependency and no image file is redistributed with it. Licence: version arXiv:2604.16041v1 of this article is distributed under the Creative Commons Attribution 4.0 International licence, as recorded in the arXiv licence metadata for this exact version and shown on the abstract page https://arxiv.org/abs/2604.16041v1. A full rights notice is delivered at licenses/arxiv-2604.16041-CC-BY-4.0.txt. Characters: every retained excerpt is pure ASCII, so the delivered engine, XeLaTeX with its default Latin Modern text font, reports no missing character.
\begin{equation} \label{def jnr}
	\begin{split}
		W(A_1, \dots, A_t) = \{  (\tr(\rho A_1),   \dots, \tr(\rho A_t))& \in \R^t:	\, \rho \in \calD \}.
	\end{split}
\end{equation}

\begin{definition}\label{defi mS}
 	Let $S$ be a subspace of $\C^n$ with $\dim(S)=r$. We define the moment $m_S$ of $S$ in a fixed orthonormal basis $\mathbb{B}=\{B_k\}_{i=1}^{\dim(\calB)}$ of Hermitian matrices of a C$^*$-subalgebra $\calB\subset M_n(\C)$, as the subset of $\R^{\dim(\calB)}$ given by
 	\begin{equation}
 		\label{def momento de un Subesp en relacion a B}
 		m_S=m_{S,\mathbb{B}}=\mathop{\bigcup}\limits_{\substack{\text{o.n. basis } \\
 				\{v_i\}_{i=1}^{r} \text{of } S}
 		}  \text{conv}\left\{ \left(\left(v_i v_i^*\right)_1,\dots, \left(v_i v_i^*\right)_{\dim(\calB)} \right)  \right\}_{i=1}^{r}
 	\end{equation}
 	where we denoted with $v_i$ the column vector, $v_i^*$ the row vector conjugated, $\left(v_i v_i^*\right)_k=\langle v_i v_i^*, B_k\rangle=\tr(v_i v_i^* B_k^*)=\tr( v_i^* B_k v_i)\in\R$ the $k^\text{th}$ coordinate of the orthogonal projection  $v_i v_i^*$ of  rank one in the basis $\mathbb{B}$ of $\calB$. See equivalent description of $m_S$ in Proposition \ref{prop 1}. 
 	
  	
\end{definition}

\begin{equation}
	\label{def matrices de densidad de S}
	\calD_{S}=\{\rho\in M^h_n(\C): \im(\rho)\subset S,\ \rho\geq 0 \text{ and } \tr(\rho)=1\}\subseteq \calD
\end{equation}

\begin{equation}\label{equ densidad escrita usando autovals y autovecs}
\rho=\sum_{i=1}^{\dim(S)} \lambda_i s_i (s_i)^*
\end{equation}

\begin{equation}\label{defi phi}
	\Phi^{\mathbb{B}}(\rho)= (\rho_1, \rho_2, \dots, \rho_{\dim(\calB)} ),
\end{equation}

	\begin{equation} \label{wi}
		w_i= \left( \tr(v_i v_i^* B_1),\tr(v_i v_i^* B_2),...,\tr(v_i v_i^* B_{\dim(\mathcal{B})})\right)\in \R^{\dim(\mathcal{B})},
	\end{equation}

\begin{proposition}[Equivalent formulas for $m_S$] \label{prop 1}
 	Let $S$ be a subspace of $\C^n$ with $\dim(S)=r>0$, $\mathbb{B}=\{B_k\}_{i=1}^{\dim(\calB)}$ be a fixed orthonormal basis of Hermitian matrices of a C$^*$-subalgebra $\calB\subset M_n(\C)$ and  $\Phi^{\mathbb{B}}:\calD_{S}\to\R^{\dim(\calB)}$ defined as in \eqref{defi phi}. Then the following statements are equivalent to the Definition  \ref{defi mS} of $m_{S,\mathbb{B}}$ ($=m_S)$.
 % 	
% 	The following statements can be used to define the moment $m_{S,\mathbb{B}}$, for $S\subsetneq\C^n$, with $\dim(S)=r$, and $\mathbb{B}=\{B_k\}_{k=1}^{\dim(\calB)}$ an orthonormal  basis of $\calB$ with $B_k^*=B_k$ for all $k$.
% 	
	\begin{enumerate}
			\item \label{item3Prop1}
		$\displaystyle{
			m_S=\mathop{\bigcup}\limits_{\substack{\text{o.n. basis}  \\
					\{v_i\}_{i=1}^{r} \text{of } S}}
			{\rm conv}\left\{ \left(\left(v_i v_i^*\right)_1,\dots, \left(v_i v_i^*\right)_{\dim(\mathbb{B})} \right)  \right\}_{i=1}^{r}
		}$
		\item \label{mS = diag DsubS} $m_S= \Phi^{\mathbb{B}}(\calD_{S})$.
		\item \label{item2Prop1} 		
		\begin{equation*}
			\begin{split}
				m_S=\text{conv}\left(\left\{(V_1,V_2,\right.\right.&\left.\left.\dots,V_{\dim(\calB)})\right.\right.\in\R^{\dim(\calB)}:\\
				&\left.\left.V=s s^*, \text{ for } s\in S, \|s\|=1  \right\} \right)
			\end{split}
		\end{equation*}
		or equivalently
		\begin{equation}\label{eq momento}
			\begin{split}
			m_S=\text{conv}\left\{\left(\langle B_1v,v\rangle, \cdots, \langle B_{dim \calB}v,v\rangle\right)
			\right. 
			& \in\R^{\dim(\calB)}: \\
			 & \left. v\in S, \|v\|=1 \right\}.		
			\end{split}
		\end{equation}
	
		\item \label{item4Prop1}
		$\displaystyle{
		m_S=\{(\tr(B_1Y),\dots,\tr(B_{\dim(\mathbb{B})} Y)): Y\in \mathcal{D}_S\} \ \ \text{(see \eqref{def matrices de densidad de S})}
			}$
		\item \label{item5Prop1}
		$
		\displaystyle{
		W(P_SB_1P_S,\dots,P_SB_{_{\dim(\calB)}}P_S)=\bigcup_{\varepsilon\in[0,1]} \varepsilon m_S}$, (\text{see} \eqref{def jnr})
		
		\item \label{item6Prop1}
		$\displaystyle{
		m_S=W(P_SB_1P_S,\dots,P_SB_{{\dim(\mathbb{B})}}P_S)\cap \left\{ x\in \mathbb{R}^{{\dim(\mathbb{B})}}: \sum_{i=1}^{\dim(\mathbb{B})} x_i=1 \right\}
		}$.
	\end{enumerate}
\end{proposition}

\begin{proof}
\begin{enumerate}

	\item[\ref{item3Prop1}.] Is \eqref{def momento de un Subesp en relacion a B} of Definition \ref{defi mS}.
	
	\item[\ref{mS = diag DsubS} $\Leftrightarrow$ \ref{item3Prop1}.] Let $\rho\in\calD_{S}$. By \eqref{equ densidad escrita usando autovals y autovecs}, there exists an orthonormal basis $\{s_i\}_{i=1}^r$ of $S$ such that
	\begin{equation} \label{eq phi de v}
		\begin{split}
		\Phi^{\mathbb{B}}(\rho)&=\sum_{i=1}^{\dim(S)} \lambda_i \Phi^{\mathbb{B}}(s_i (s_i)^*)\\
		&=\sum_{i=1}^{\dim(S)} \lambda_i \left( (s_i (s_i)^*)_1, (s_i (s_i)^*)_2, \dots, (s_{i} (s_{i})^*)_{\dim(\calB)}\right), 
	\end{split}
	\end{equation}
	with $\sum_{i=1}^{\dim(S)}\lambda_i=1$ and $\lambda_i\geq 0$. Observe that, since $s_i s_i^*\geq 0$, $\tr(s_i s_i^*)=1$ and
	\begin{equation*}
		\begin{split}
	\Phi^{\mathbb{B}}(s_i (s_i)^*)&=\left(\tr(s_is_i^*B_1), \tr(s_is_i^*B_2),..., \tr(s_is_i^*B_{\dim(\mathcal{B})}) \right)\\
	&=\left(\langle B_1s_i,s_i\rangle,\langle B_2s_i,s_i\rangle,...,\langle B_{\dim(\mathcal{B})}s_i,s_i\rangle \right),		
		\end{split}
	\end{equation*}
	then each $\Phi^{\mathbb{B}}(s_i (s_i)^*)\in W\left(\{B_k\}_{k=1}^{\dim(\mathcal{B})} \right)$.
	Then, $\Phi^{\mathbb{B}}(\rho)\in\text{conv}\{\Phi^{\mathbb{B}}(s_i (s_i)^*)\}_{i=1}^{\dim(\mathcal{B})}\Rightarrow \Phi^{\mathbb{B}}(\rho)\in m_S$ which implies that $\Phi^{\mathbb{B}}(\calD_{S})\subseteq m_S$.
		
	In order to prove the other inclusion, take $x\in m_S$. Then there exist  $\{w_i\}_{i=1}^{r}$ as in \eqref{wi} such that
	 $x=\sum_{i=1}^{r}\alpha_iw_i$, with $0\leq \alpha_i$ for all $i$ and $\sum_{i=1}^{r}\alpha_i=1$. Then, 
	 \begin{align*}
	 	x&= \sum_{i=1}^{r}\alpha_i\left( \tr(v_i v_i^* B_1),\tr(v_i v_i^* B_2),...,\tr(v_i v_i^* B_{\dim(\mathcal{B})})\right) \\
	 	&=  \left( \sum_{i=1}^{r}\alpha_i\tr(v_i v_i^* B_1),\sum_{i=1}^{r}\alpha_i\tr(v_i v_i^* B_2),...,\sum_{i=1}^{r}\alpha_i\tr(v_i v_i^* B_{\dim(\mathcal{B})})\right) \\
	 	&= \left(\tr\left( \left( \sum_{i=1}^{r}\alpha_iv_i v_i^*\right)  B_1\right) 
	 	%,\tr\left(\left(  \sum_{i=1}^{r}\alpha_iv_i v_i^*\right)  B_2\right) 
	 	,...,\tr\left( \left( \sum_{i=1}^{r}\alpha_iv_i v_i^*\right)  B_{\dim(\mathcal{B})}\right) \right) \\
	 	&= \left(\tr\left( \rho B_1\right) ,\tr\left(\rho  B_2\right) ,...,\tr\left( \rho B_{\dim(\mathcal{B})}\right) \right)
	 \end{align*}
	 for $\rho=\sum_{i=1}^{r}\alpha_iv_i v_i^*$.	 
	 Since $\{v_i\}_{i=1}^r\subset S$ is an orthonormal basis of $S$, follows that $\rho\geq 0$, 
	 $\tr(\rho)=1$ and then $\rho \in \mathcal{D}_S$. Therefore $x= \Phi(\rho)\in \Phi(\mathcal{D}_S)$.	 
	 
	\item[\ref{item2Prop1} $\Leftrightarrow$ \ref{item3Prop1}.] Let $x\in m_S$. By the previous item there exist  $\rho\in\calD_{S}$ such that $x=\Phi(\rho)$. Observing \eqref{eq phi de v}, we deduce that 
	\begin{equation*}
		\begin{split}
	\Phi^\mathbb{B}(\rho)\in \text{conv}  \left\{ (V_1,V_2,\dots,V_{\dim(\calB)})\in\R^{\dim(\calB)}:   
	V=s s^*, \text{ for }  
	s\in S,    
	  	    \|s\|=1  \right\}  ,		
		\end{split}
	\end{equation*}
	where $V_k=\langle V, B_k\rangle_{tr}=\tr(V B_k)$ with $\mathbb{B}=\{B_k\}_{k=1}^{\dim \calB}$ the fixed basis of $\calB$.
	
	Now we prove the reverse inclusion. Let 
	\[y\in \text{conv} \left(\left\{(V_1, \dots,V_{\dim(\calB)})\in\R^{\dim(\calB)}:V=s s^*, \text{ for } s\in S, \|s\|=1  \right\} \right).\] 
	There exist $s_j\in S$, $\|s_j\|=1$ such that 
	$	y = \sum_{j=1}^{k}\beta_js_js_j^*$,
	with $0\leq \beta_j$ for all $1\leq j\leq k$ and $\sum_{j=1}^{k}\beta_j=1$. Then,
	\begin{align*}
		y& = \sum_{j=1}^{k}\beta_j\left( (s_js_j^*)_1,(s_js_j^*)_2,..., (s_js_j^*)_{\dim(\calB)}\right) \\
		& = \sum_{j=1}^{k}\beta_j\left( \tr(s_js_j^*B_1),\tr(s_js_j^*B_2),..., \tr(s_js_j^*B_{\dim(\calB)})\right),\\
%		& = \left(\sum_{j=1}^{k}\beta_j \tr(s_js_j^*B_1),\sum_{j=1}^{k}\beta_j\tr(s_js_j^*B_2),..., \sum_{j=1}^{k}\beta_j\tr(s_js_j^*B_{\dim(\calB)})\right) \\
	\end{align*}
		where $\left( \tr(s_js_j^*B_1),\tr(s_js_j^*B_2),..., \tr(s_js_j^*B_{\dim(\calB)})\right)\in m_S$ for every $1\leq j\leq k$. Therefore, $y\in m_S$.
	\end{enumerate}
 
	\begin{enumerate}
		\item[\ref{item4Prop1} $\Rightarrow$ \ref{item3Prop1}.]  If $Y\in \mathcal{D}_S$, then $Y=\sum_{i=1}^r \sigma_i\, y_iy_i^*$, with $\sum_{i=1}^r\sigma_i=1$, $\sigma_i\geq 0$, and $y_i\in S$ with $y_i\perp y_j$ for $i\neq j$ ($\{y_i\}_{i=1}^r$ is an orthonormal basis of $S$). Then 
	\begin{equation*}\begin{split}
			(\tr(B_1 Y),\dots,\tr(B_{\dim(\mathbb{B})} Y))&=\sum_{i=1}^r\sigma_i \left(\tr(B_1y_iy_i^*),\dots ,\tr(B_{\dim(\mathbb{B})} y_iy_i^*)\right)\\
			&=	\sum_{i=1}^r\sigma_i \left((y_iy_i^*)_1,\dots ,(y_iy_i^*)_{\dim(\mathbb{B})}\right)
			%\in \bigcup\limits_{\substack{\text{o.n. basis}  \\				\{v_i\}_{i=1}^{r} \subset S}} \text{conv}
		\end{split}
	\end{equation*}
	which belongs to $\mathop{\bigcup}\limits_{\substack{\text{o.n. basis}  \\
			\{v_i\}_{i=1}^{r} \text{of } S}}
	\text{conv}\left\{ \left(\left(v_i v_i^*\right)_1,\dots, \left(v_i v_i^*\right)_{{\dim(\mathbb{B})}} \right)  \right\}_{i=1}^{r}$.
	
	\item[\ref{item3Prop1} $\Rightarrow$ \ref{item4Prop1}.] Given a convex combination $ \sum_{i=1}^r \sigma_i\ \left(\tr(B_1 y_iy_i^*),\dots,\tr(B_{\dim(\mathbb{B})} y_iy_i^*) \right)$, where $\{y_i\}_{i=1}^r$ is  an orthonormal basis of $S$, $\sum_{i=1}^r\sigma_i=1$ with $\sigma_i\geq 0$, it is easy to see that $Y=\sum_{i=1}^r \sigma_i \, y_iy_i^*\in\DS$.  
 
	
	\item[\ref{item4Prop1} $\Rightarrow$ \ref{item5Prop1}.] 
	Recall that
	\begin{equation*}
		\begin{split}
			W&(P_SB_1P_S,\dots,P_SB_{{\dim(\mathbb{B})}}P_S)=\\
			&=\{\left(\tr(\rho P_SB_1P_S), \dots, \tr(\rho P_SB_{\dim(\mathbb{B})}P_S)\right)\in\R^{\dim(\mathbb{B})}:
			 \\&  
			\qquad \rho\in M_n(\C),\rho\geq 0, \tr(\rho)=1\}.
		\end{split}
	\end{equation*} 
	Also observe that choosing $\rho=\frac1{\dim(S^\perp)}P_{S^\perp}$ follows that $(0,\dots,0)\in W(P_SB_1P_S,\dots,P_SB_{{\dim(\mathbb{B})}}P_S)$. Now for $\rho\in\mathcal{D}$, if $\tr(P_S\rho P_S)\neq 0$ we can write
	\begin{equation*}
		\begin{split}
			(\tr(\rho P_SB_1 P_S),\dots,&\tr(\rho P_SB_{\dim(\mathbb{B})} P_S))=\\
			&=(\tr(P_S\rho P_SB_1 ),\dots,\tr(P_S\rho P_SB_{\dim(\mathbb{B})}))\\
			&=
			\tr(P_S\rho P_S)(\tr(\rho_S B_1 ),\dots,\tr(\rho_S B_{\dim(\mathbb{B})}))		
		\end{split}
	\end{equation*}
	where $\rho_S=\frac1{\tr(P_S\rho P_S)}P_S\rho P_S\in \calD_S$. Hence, since $\tr(P_S\rho P_S)\leq 1$, we proved that every element of $W(P_SB_1P_S,\dots,P_SB_{{\dim(\mathbb{B})}}P_S)$ can be written as $\varepsilon(\tr(\rho_S B_1 ),\dots,\tr(\rho_S B_{\dim(\mathbb{B})}))$ with $\varepsilon\in[0,1]$ and $\rho_S\in \calD_S$. Then, since $(\tr(\rho_S B_1 ),\dots,\tr(\rho_S B_{\dim(\mathbb{B})}))\in m_S$, we have proved the inclusion $\subset$ of the equality in \ref{item5Prop1}.
	
	Now we will consider the inclusion $\supset$ of \ref{item5Prop1}. Suppose $C=\varepsilon m$, with $\varepsilon\in[0,1]$ and $m\in m_S$. Then $C=\varepsilon (\tr(B_1 Y),\dots ,\tr(B_{\dim(\mathbb{B})} Y))$ with $Y\in\calD_S$. The case when $\varepsilon=0$ is trivial. Now if $\varepsilon\in(0,1]$, choose $Z= \frac{1-\varepsilon}{\varepsilon} \frac1{\dim S^\perp}P_{S^\perp}$ and consider $\rho=\varepsilon(Y+Z)$. Then $\tr(\rho)=\varepsilon (\tr(Y)+\tr(Z))=\varepsilon (1+\frac{1-\varepsilon}{\varepsilon})=1$, $\rho\geq 0$ which implies that $\rho\in\calD$. Also observe that
	\begin{equation*}
		\begin{split}
			\varepsilon(\tr(B_1 Y),\dots ,\tr(B_{\dim(\mathbb{B})} Y))&=
			\varepsilon(\tr(B_1 P_S Y P_S),\dots ,\tr(B_{\dim(\mathbb{B})} P_SY P_S)) \\
			&=\varepsilon
			(\tr(B_1 P_S (Y+Z) P_S),\dots ,\\
			& \hspace{3cm} \tr(B_{\dim(\mathbb{B})} P_S(Y+Z) P_S))\\
			&=(\tr(B_1 P_S \rho P_S),\dots ,\tr(B_{\dim(\mathbb{B})} P_S\rho P_S))
			\end{split}
	\end{equation*}
	which belongs to $W(P_SB_1P_S,\dots,P_SB_{{\dim(\mathbb{B})}}P_S)$.
	
	
	\item[\ref{item5Prop1} $\Rightarrow$ \ref{item4Prop1}.] This implication follows considering $\varepsilon=1$ and $\rho=Y\in\calD_S$ in \ref{item5Prop1}.
	
	\item[\ref{item6Prop1}.] This statement can be proved using items \ref{mS = diag DsubS} and \ref{item5Prop1} of this proposition.
\end{enumerate}
\end{proof}

\end{document}
